\chapter{Future Work}

\section{Integration of Deep Learning Models}

The Gradient Boosting model developed in this thesis proved to be highly effective for its intended purpose of rapid, structure-agnostic screening. However, the field of machine learning is constantly evolving, and the rise of deep learning presents exciting opportunities to build upon this work and overcome some of its inherent limitations. Future research should focus on the integration of more advanced deep learning architectures, particularly those designed to work with structured data like crystal lattices.

A promising and highly relevant direction is the application of \textbf{Graph Neural Networks (GNNs)}. Crystal structures can be naturally represented as graphs, where atoms are the nodes and the chemical bonds between them are the edges. GNNs are specifically designed to operate on such graph-structured data and have shown state-of-the-art performance in materials science applications \cite{Xie2018}. A key advantage of this approach is its ability to learn feature representations directly from the crystal graph, a process known as end-to-end learning. This would bypass the need for manual feature engineering and could potentially allow the model to discover more subtle and complex structure-property relationships than our current composition-based features can capture. Architectures like the Crystal Graph Convolutional Neural Network (CGCNN) could be trained to predict ionic conductivity with even greater accuracy by explicitly considering the local atomic environment and connectivity of the crystal lattice \cite{Xie2018}.

Furthermore, deep learning models are exceptionally well-suited for \textbf{multi-target prediction}. A significant limitation of the current work is its focus on a single property (ionic conductivity). A future GNN-based model could be designed as a multi-task framework to predict several critical properties simultaneously. For instance, the model could be trained to predict not only ionic conductivity but also thermodynamic stability (e.g., energy above the convex hull from the Materials Project) and the electrochemical stability window \cite{Jain2013}. Such a multi-property model would provide a much more holistic evaluation of a candidate material, leading to a more efficient and targeted discovery pipeline.

Finally, one could explore the use of \textbf{generative deep learning models} for inverse design. While predictive models screen existing candidates, generative models, such as Generative Adversarial Networks (GANs) or Variational Autoencoders (VAEs), can be trained to generate entirely new, hypothetical crystal structures that are optimized for high ionic conductivity and stability \cite{Kim2020}. This represents a paradigm shift from merely accelerating the screening process to actively designing novel materials from the ground up, and it stands as a long-term goal for the field.

\section{Exploring New Feature Sets}

The composition-based feature set developed in this thesis proved highly effective for rapid screening. However, its predictive power is ultimately limited by its structure-agnostic design. A significant avenue for future work lies in enhancing the model by developing and incorporating more sophisticated, physically-motivated features that can provide a proxy for structural information without incurring the high computational cost of full first-principles calculations.

One promising approach is to leverage features derived from \textbf{Voronoi tessellation}. Given a relaxed crystal structure (which can often be obtained from databases or low-cost simulations), a Voronoi analysis can be performed to partition the structure into polyhedra around each atom. From this analysis, a rich set of descriptors can be extracted, such as the volume and shape of atomic polyhedra, coordination numbers, and measures of local packing density \cite{Ward2017}. These features provide a much more detailed description of the local atomic environment than simple compositional averages and could help the model better distinguish between materials with similar compositions but different local structures.

Another powerful and computationally inexpensive technique to explore is the \textbf{Bond Valence Sum (BVS)} method. The BVS model is an empirical model that can be used to validate crystal structures and predict the likely pathways for ion migration by identifying sites with low valence mismatch for the mobile ion \cite{Adams2001}. Features derived from BVS analysis, such as the bond valence energy landscape, could provide direct, quantitative information about the favorability of ion hopping within a structure, making them highly potent predictors for ionic conductivity.

Finally, future work could focus on automated feature engineering techniques. Instead of manually crafting features based on chemical intuition, methods like genetic programming or symbolic regression could be employed to automatically discover complex mathematical combinations of primary elemental properties that have a strong, non-obvious correlation with ionic conductivity. This data-driven approach to feature creation could uncover novel descriptors that lead to even more accurate and insightful models. By systematically exploring these richer feature sets, future iterations of this work can bridge the gap between simple compositional models and computationally intensive, structure-based deep learning methods, leading to predictive tools with both enhanced accuracy and interpretability.
